% questions.tex —— 2009 年·普通高等学校招生考试·北京卷·理科 · 题目内容片段（无答案）
% 由 tools/generate_projects.py 依来源总账自 2009/2009_数学真题_北京理_试题.md 生成，勿手改（改源 MD 后重跑）


状态：needs-review；依固定候选 PDF 两页（印刷页 499–500）逐题核验，图形已矢量重绘，正式构建与逐页验收见回执。源未标分值，不补分值。

\gkpart{一、选择题}

\begin{question}
\gkti{1}在复平面内，复数 $z=\mathrm i(1+2\mathrm i)$ 对应的点位于（　）

（A）第一象限

（B）第二象限

（C）第三象限

（D）第四象限

\end{question}
\begin{question}
\gkti{2}已知向量 $\boldsymbol a$、$\boldsymbol b$ 不共线，$\boldsymbol c=k\boldsymbol a+\boldsymbol b$（$k\in\mathbf R$），$\boldsymbol d=\boldsymbol a-\boldsymbol b$．如果 $\boldsymbol c\parallel\boldsymbol d$，那么（　）

（A）$k=1$ 且 $\boldsymbol c$ 与 $\boldsymbol d$ 同向

（B）$k=1$ 且 $\boldsymbol c$ 与 $\boldsymbol d$ 反向

（C）$k=-1$ 且 $\boldsymbol c$ 与 $\boldsymbol d$ 同向

（D）$k=-1$ 且 $\boldsymbol c$ 与 $\boldsymbol d$ 反向

\end{question}
\begin{question}
\gkti{3}为了得到函数 $y=\lg\dfrac{x+3}{10}$ 的图象，只需把函数 $y=\lg x$ 的图象上所有的点（　）

（A）向左平移 $3$ 个单位长度，再向上平移 $1$ 个单位长度

（B）向右平移 $3$ 个单位长度，再向上平移 $1$ 个单位长度

（C）向左平移 $3$ 个单位长度，再向下平移 $1$ 个单位长度

（D）向右平移 $3$ 个单位长度，再向下平移 $1$ 个单位长度

\end{question}
\begin{question}
\gkti{4}若正四棱柱 $ABCD-A_1B_1C_1D_1$ 的底面边长为 $1$，$AB_1$ 与底面 $ABCD$ 成 $60^\circ$ 角，则 $A_1C_1$ 到底面 $ABCD$ 的距离为（　）

（A）$\frac{\sqrt3}{3}$

（B）$1$

（C）$\sqrt2$

（D）$\sqrt3$

\end{question}
\begin{question}
\gkti{5}「$\alpha=\dfrac{\pi}{6}+2k\pi$（$k\in\mathbf Z$）」是「$\cos2\alpha=\dfrac12$」的（　）

（A）充分而不必要条件

（B）必要而不充分条件

（C）充分必要条件

（D）既不充分也不必要条件

\end{question}
\begin{question}
\gkti{6}若 $(1+\sqrt2)^5=a+b\sqrt2$（$a,b$ 为有理数），则 $a+b=$（　）

（A）$45$

（B）$55$

（C）$70$

（D）$80$

\end{question}
\begin{question}
\gkti{7}用 $0$ 到 $9$ 这 $10$ 个数字，可以组成没有重复数字的三位偶数的个数为（　）

（A）$324$

（B）$328$

（C）$360$

（D）$648$

\end{question}
\begin{question}
\gkti{8}点 $P$ 在直线 $l:y=x-1$ 上，若存在过 $P$ 的直线交抛物线 $y=x^2$ 于 $A,B$ 两点，且 $|PA|=|AB|$，则称点 $P$ 为「$\mathcal A$ 点」，那么下列结论中正确的是（　）

（A）直线 $l$ 上的所有点都是「$\mathcal A$ 点」

（B）直线 $l$ 上仅有有限个点是「$\mathcal A$ 点」

（C）直线 $l$ 上的所有点都不是「$\mathcal A$ 点」

（D）直线 $l$ 上有无穷多个点（但不是所有的点）是「$\mathcal A$ 点」

\end{question}
\gkpart{二、填空题}

\begin{question}
\gkti{9}$\lim\limits_{x\to1}\dfrac{x\sqrt x-x}{x-1}=$\gkblank{}．

\end{question}
\begin{question}
\gkti{10}若实数 $x,y$ 满足 $\begin{cases}x+y-2\geqslant0,\\ x\leqslant4,\\ y\leqslant5,\end{cases}$ 则 $s=y-x$ 的最小值为\gkblank{}．

\end{question}
\begin{question}
\gkti{11}设 $f(x)$ 是偶函数．若曲线 $y=f(x)$ 在点 $(1,f(1))$ 处的切线的斜率为 $1$，则该曲线在 $(-1,f(-1))$ 处的切线的斜率为\gkblank{}．

\end{question}
\begin{question}
\gkti{12}椭圆 $\dfrac{x^2}{9}+\dfrac{y^2}{2}=1$ 的焦点为 $F_1,F_2$，点 $P$ 在椭圆上．若 $|PF_1|=4$，则 $|PF_2|=$\gkblank{}；$\angle F_1PF_2$ 的大小为\gkblank{}．

\end{question}
\begin{question}
\gkti{13}若函数 $f(x)=\begin{cases}\dfrac1x,&x<0,\\[4pt] \left(\dfrac13\right)^x,&x\geqslant0,\end{cases}$ 则不等式 $|f(x)|\geqslant\dfrac13$ 的解集为\gkblank{}．

\end{question}
\begin{question}
\gkti{14}已知数列 $\{a_n\}$ 满足：$a_{4n-3}=1$，$a_{4n-1}=0$，$a_{2n}=a_n$（$n\in\mathbf N^*$），则 $a_{2009}=$\gkblank{}；$a_{2014}=$\gkblank{}．

\end{question}
\gkpart{三、解答题}

\begin{question}
\gkti{15}在 $\triangle ABC$ 中，角 $A,B,C$ 的对边分别为 $a,b,c$，$B=\dfrac{\pi}{3}$，$\cos A=\dfrac45$，$b=\sqrt3$．

（1）求 $\sin C$ 的值；

（2）求 $\triangle ABC$ 的面积．

\Needspace{0.4\textheight}
\end{question}
\begin{question}
\gkti{16}如图，在三棱锥 $P-ABC$ 中，$PA\perp$ 底面 $ABC$，$PA=AB$，$\angle ABC=60^\circ$，$\angle BCA=90^\circ$，点 $D,E$ 分别在棱 $PB,PC$ 上，且 $DE\parallel BC$．
\nopagebreak[4]
\begin{center}% 2009 年北京卷理科第 16 题：三棱锥 P-ABC，PA⊥底面，D、E 在 PB、PC 上且 DE∥BC。
% 来源 PDF：普通高考/2009/2009北京理.pdf，SHA256 f5ca2006abecc7e9529b7adcddcda172d995b4910e43445e33125234efcfc4c2
% 原页：第 1 页（A3 三栏）第 3 栏中部，印刷页 499。
% 重绘依据：P 在顶点，A 在左、B 在右、C 在下方居前；PA、PB、PC、AB、AC、BC 实线，
%   D 在 PB 上、E 在 PC 上，DE 水平且平行 BC；标注按原图。
\begin{tikzpicture}[scale=1.15, semithick, every node/.style={inner sep=1.5pt}]
  \coordinate (P) at (0.7,2.9);
  \coordinate (A) at (-1.4,1.4);
  \coordinate (B) at (2.2,1.4);
  \coordinate (C) at (0.4,0.2);
  \coordinate (D) at ($(P)!0.55!(B)$);
  \coordinate (E) at ($(P)!0.55!(C)$);
  \draw (P) -- (A);
  \draw (P) -- (B);
  \draw (P) -- (C);
  \draw (A) -- (B);
  \draw (A) -- (C);
  \draw (C) -- (B);
  \draw (D) -- (E);
  \fill (D) circle (1.2pt);
  \fill (E) circle (1.2pt);
  \node[anchor=south] at (P) {$P$};
  \node[anchor=east] at (A) {$A$};
  \node[anchor=west] at (B) {$B$};
  \node[anchor=north] at (C) {$C$};
  \node[anchor=west] at (D) {$D$};
  \node[anchor=east] at (E) {$E$};
\end{tikzpicture}
\end{center}
\nopagebreak[4]
（1）求证：$BC\perp$ 平面 $PAC$；

（2）当 $D$ 为 $PB$ 的中点时，求 $AD$ 与平面 $PAC$ 所成的角的大小；

（3）是否存在点 $E$ 使得二面角 $A-DE-P$ 为直二面角？并说明理由．

\end{question}
\begin{question}
\gkti{17}某学生在上学路上要经过 $4$ 个路口，假设在各路口是否遇到红灯是相互独立的，遇到红灯的概率都是 $\dfrac13$，遇到红灯时停留的时间都是 $2$ min．

（1）求这名学生在上学路上到第三个路口时首次遇到红灯的概率；

（2）求这名学生在上学路上因遇到红灯停留的总时间 $\xi$ 的分布列及期望．

\end{question}
\begin{question}
\gkti{18}设函数 $f(x)=x\mathrm e^{kx}$（$k\neq0$）．

（1）求曲线 $y=f(x)$ 在点 $(0,f(0))$ 处的切线方程；

（2）求函数 $f(x)$ 的单调区间；

（3）若函数 $f(x)$ 在区间 $(-1,1)$ 内单调递增，求 $k$ 的取值范围．

\end{question}
\begin{question}
\gkti{19}已知双曲线 $C:\dfrac{x^2}{a^2}-\dfrac{y^2}{b^2}=1$（$a>0,b>0$）的离心率为 $\sqrt3$，右准线方程为 $x=\dfrac{\sqrt3}{3}$．

（1）求双曲线 $C$ 的方程；

（2）设直线 $l$ 是圆 $O:x^2+y^2=2$ 上动点 $P(x_0,y_0)$（$x_0y_0\neq0$）处的切线，$l$ 与双曲线 $C$ 交于不同的两点 $A,B$，证明 $\angle AOB$ 的大小为定值．

\end{question}
\begin{question}
\gkti{20}已知数集 $A=\{a_1,a_2,\cdots,a_n\}$（$1\leqslant a_1<a_2<\cdots<a_n$，$n\geqslant2$）具有性质 $P$：对任意的 $i,j$（$1\leqslant i\leqslant j\leqslant n$），$a_ia_j$ 与 $\dfrac{a_j}{a_i}$ 两数中至少有一个属于 $A$．

（1）分别判断数集 $\{1,3,4\}$ 与 $\{1,2,3,6\}$ 是否具有性质 $P$，并说明理由；

（2）证明：$a_1=1$，且 $\dfrac{a_1+a_2+\cdots+a_n}{a_1^{-1}+a_2^{-1}+\cdots+a_n^{-1}}=a_n$；

（3）证明：当 $n=5$ 时，$a_1,a_2,a_3,a_4,a_5$ 成等比数列．

\end{question}